% Book-local P03.1 affine ground-field curve-test proof,
% 12 September 2026.  No novelty or priority claim is made.

\begin{proof}[\Cref{lem:p031-open-pullback-usc}]
For $x\in X$, write the upper-semicontinuous regularization as the
neighborhood-filter limsup
\[
                  a^*(x)=\inf_{U\ni x}\sup_{z\in U}a(z).
\]
Fix $y\in Y$ and put $x=q(y)$.  If $N$ is a neighborhood of $y$,
then $q(N)$ is a neighborhood of $x$ because $q$ is open.
Conversely, if $U$ is a neighborhood of $x$, then $q^{-1}(U)$ is a
neighborhood of $y$ and, by surjectivity, $q(q^{-1}(U))=U$.  Hence the
sets $q(N)$, as $N$ runs over the neighborhoods of $y$, are cofinal
among the neighborhoods of $x$.  Therefore
\[
 \begin{split}
 (a\circ q)^*(y)
   &=\inf_{N\ni y}\sup_{z\in N}a(q(z))\\
   &=\inf_{N\ni y}\sup_{x'\in q(N)}a(x')
    =a^*(x),
 \end{split}
\]
which is the claimed identity.
\end{proof}

\begin{proof}[
\Cref{thm:p031-affine-ground-field-perron}]
Put $s:=\sup_i u_i$ and $u:=s^*$.  Local upper boundedness makes $u$
an upper-semicontinuous function with values in
$\mathbf R\cup\{-\infty\}$.

Let $C$ be a smooth connected strict $K$-analytic curve and let
$g:C\to V$ be a morphism.  We prove that $u\circ g$ is subharmonic.  It
is enough, by \cite[Proposition~3.5(3)]{GR25Classical}, to work on each
connected component of $\operatorname{Int}(C)$.  If $u\circ g$ is
identically $-\infty$ there, there is nothing to prove.  Otherwise let
$D$ be a strict affinoid domain in that component and let
$h:D\to\mathbf R$ be harmonic with
\begin{equation}
\label{eq:p031-ground-field-boundary-hypothesis}
                         u\circ g\leq h
                         \qquad\text{on }\partial D.
\end{equation}
We must prove the same inequality on $D$.

We first construct a small analytic tube around the graph of $g|_D$.
Write $z_1,\ldots,z_n$ for the affine coordinates and
$g_j:=z_j\circ g\in\mathcal O(D)$.  On
$D\times\mathbf A_K^{n,\mathrm{an}}$, the translation
\[
  \Psi(t,w_1,\ldots,w_n)
     =(t,g_1(t)+w_1,\ldots,g_n(t)+w_n)
\]
is an analytic automorphism.  The inverse image
$\Psi^{-1}(D\times V)$ is an ordinary open neighborhood of the zero
section.  Choose a sequence $\rho_m\in|K^\times|$ decreasing to zero.
Compactness shows that, for all sufficiently large $m$,
\begin{equation}
\label{eq:p031-ground-field-tube-inclusion}
 D\times\mathbb B^n_{\rho_m}
       \subseteq\Psi^{-1}(D\times V).
\end{equation}
Indeed, otherwise the nonempty nested compact sets
\[
 \bigl(D\times\mathbb B^n_{\rho_m}\bigr)
        \setminus\Psi^{-1}(D\times V)
\]
would have nonempty intersection, whereas their intersection is contained
in the zero section and the zero section lies in
$\Psi^{-1}(D\times V)$.

Fix a radius $\rho$ satisfying
\eqref{eq:p031-ground-field-tube-inclusion} and set
\[
 T_\rho:=D\times\mathbb B^n_\rho,
 \qquad
 p_\rho(t,w):=g(t)+w\in V,
 \qquad
 H_\rho(t,w):=h(t).
\]
Thus $T_\rho$ is a $K$-affinoid space.  Let
$\sigma_\rho:D\to T_\rho$ be the relative Gauss section.  We need two
elementary properties of this tube.

First,
\begin{equation}
\label{eq:p031-relative-gauss-boundary}
                  \sigma_\rho(\Gamma(D))
                  \quad\text{is a boundary of }T_\rho,
\end{equation}
where $\Gamma(D)=\partial D$ is the finite Shilov boundary of the strict
affinoid curve $D$, as recalled in
\cite[\S3.1(3)]{GR25Classical}.  To see this, write an analytic function on
the tube as $F=\sum_{\alpha}a_\alpha w^\alpha$, with
$a_\alpha\in\mathcal O(D)$.  The defining property of the relative
Gauss point gives
\[
 |F(\sigma_\rho(x))|
        =\max_\alpha |a_\alpha(x)|\rho^{|\alpha|}.
\]
For every point of the fiber over $x$, the value of $|F|$ is bounded
above by the right side.  Since every coefficient $a_\alpha$ attains its
supremum on $\Gamma(D)$, and the weighted coefficient norms tend to
zero, we get
\[
 \sup_{T_\rho}|F|
   =\max_{b\in\Gamma(D)}|F(\sigma_\rho(b))|.
\]
This identifies the displayed finite set as a boundary of $T_\rho$.
Consequently the Shilov boundary of $T_\rho$ is contained in that set.

Second, if $t_0\in D\setminus\partial D$, then $p_\rho$ is open on an
ordinary neighborhood of $(t_0,0)$.  In the $w$-coordinates, put
\[
 U_\rho^\circ
    :=\operatorname{Int}(D)\times\mathbb B_{\rho}^{n,\circ}
       \subseteq T_\rho,
\]
where the second factor is the open polydisc.  Before imposing the closed
tube, $p_\rho$ is, through the translation $\Psi$, the projection
\[
             D\times\mathbf A_K^{n,\mathrm{an}}
                     \longrightarrow\mathbf A_K^{n,\mathrm{an}}.
\]
The structural map of the smooth curve $\operatorname{Int}(D)$ is
quasi-smooth and boundaryless.  Base change, the translation $\Psi$, and
restriction to the ordinary open set $U_\rho^\circ$ preserve these two
properties.  Thus
\[
                   p_\rho|_{U_\rho^\circ}:U_\rho^\circ\longrightarrow V
\]
is quasi-smooth and boundaryless, and is open by
\cite[Corollary~6.2.5]{Ducros18Families}.  Since
$(t_0,0)\in U_\rho^\circ$, this proves the local assertion.

Return to the boundary comparison.  As $\rho\to0$, for every
$b\in\Gamma(D)$ one has
\begin{equation}
\label{eq:p031-gauss-convergence}
                 p_\rho(\sigma_\rho(b))\longrightarrow g(b).
\end{equation}
This follows directly from the Gauss formula: on every fixed large tube,
the value of a convergent power series at the relative Gauss point tends
to the value of its constant coefficient at the zero section.  The set
$\Gamma(D)$ is finite.  Upper semicontinuity of $u$, together with
\eqref{eq:p031-ground-field-boundary-hypothesis}, therefore shows that for every
$\varepsilon>0$ one may choose $\rho\in|K^\times|$ small enough that
\begin{equation}
\label{eq:p031-gauss-bound}
 u(p_\rho(\sigma_\rho(b)))
       \leq h(b)+\varepsilon
       \qquad(b\in\Gamma(D)).
\end{equation}

The harmonic function $h$ is pluriharmonic on $D$: this is
\cite[Definition~3.13 and \S4.1.2]{GR25Classical}.  Hence $H_\rho$ is
pluriharmonic by analytic pullback, using
\cite[Proposition~7.2(7)]{GR25Classical}.  For every $i$, the function
\[
                         u_i\circ p_\rho-H_\rho
\]
is consequently classically psh on $T_\rho$.  Since
$u_i\leq s\leq u$, equation \eqref{eq:p031-gauss-bound} bounds this
function by $\varepsilon$ on $\sigma_\rho(\Gamma(D))$.  The global
affinoid maximum principle \cite[Theorem~6.4]{GR25Classical}, together
with \eqref{eq:p031-relative-gauss-boundary}, gives
\[
                    u_i\circ p_\rho-H_\rho\leq\varepsilon
                    \qquad\text{on }T_\rho
\]
for every $i$, and hence
\begin{equation}
\label{eq:p031-raw-tube-bound}
                    s\circ p_\rho\leq H_\rho+\varepsilon
                    \qquad\text{on }T_\rho.
\end{equation}

Fix $t_0\in D\setminus\partial D$.  By the local openness proved above,
we may choose an ordinary neighborhood $N$ of $(t_0,0)$ in $T_\rho$
such that $p_\rho(N)$ is an ordinary open neighborhood of $g(t_0)$.
Shrinking $N$, continuity of $h$ also gives
$H_\rho<h(t_0)+\varepsilon$ on $N$.  It follows from
\eqref{eq:p031-raw-tube-bound} that
\[
                    s<h(t_0)+2\varepsilon
                    \qquad\text{on }p_\rho(N).
\]
The definition of usc regularization now yields
\[
                         u(g(t_0))\leq h(t_0)+2\varepsilon.
\]
Letting $\varepsilon\downarrow0$ proves $u\circ g\leq h$ at every
point of $D\setminus\partial D$; the boundary inequality was assumed in
\eqref{eq:p031-ground-field-boundary-hypothesis}.  Thuillier's comparison criterion
shows that $u\circ g$ is subharmonic on each component of
$\operatorname{Int}(C)$, and
\cite[Proposition~3.5(3)]{GR25Classical} extends it across the relative
boundary of $C$.  This is exactly the property in
\Cref{def:p031-ground-field-curvewise-psh}.
\end{proof}

\begin{proof}[
\Cref{cor:p031-scalar-compatibility-upgrade}]
Fix an analytic extension $K'/K$.  Pointwise supremum commutes with
pullback, so
\[
             s\circ\pi_{K'/K}=\sup_i(u_i\circ\pi_{K'/K}).
\]
The pulled-back functions are classically psh by
\cite[Proposition~4.2(5)]{GR25Classical} and remain locally upper bounded.
Applying \Cref{thm:p031-affine-ground-field-perron} over $K'$ shows
that $(s\circ\pi_{K'/K})^*$ is $K'$-curvewise psh.  Under
\eqref{eq:p031-scalar-compatibility}, this function equals
$u\circ\pi_{K'/K}$.  Since this holds for every $K'/K$, the smooth
strict curve test in \cite[Definition~4.1 and \S4.1.3]{GR25Classical}
shows that $u$ is classically psh.

If $\pi_{K'/K}$ is open, it is also continuous and surjective; the latter
two properties of scalar projection are recalled in
\cite[\S2.9]{GR25Classical}.  Apply
\Cref{lem:p031-open-pullback-usc} to $a=s$ and use $u=s^*$ to get
\[
 u\circ\pi_{K'/K}
   =s^*\circ\pi_{K'/K}
   =(s\circ\pi_{K'/K})^*,
\]
which is the desired identity for that extension.
\end{proof}

\begin{proof}[
\Cref{cor:p031-affine-obstacle-ground-field}]
A locally Fubini--Study function is classically psh by the logarithm and
sheaf properties, together with closure under positive linear combinations
and finite maxima; see
\cite[Proposition~4.2(1)--(3)]{GR25Classical}.  A semipositive real-PL
function is classically psh by \cite[Theorem~4.3]{GR25Classical}.  The
upper-semicontinuous obstacle $f$, which has no $+\infty$ value, is
locally bounded above, and therefore bounds the whole family
$\mathcal A_f$ locally.  If this family is nonempty, apply
\Cref{thm:p031-affine-ground-field-perron}; the scalar-compatible
conclusion follows from
\Cref{cor:p031-scalar-compatibility-upgrade}.  If the family is empty,
then $P_f=-\infty$, which is already classically psh by
\cite[Definitions~3.4 and~4.1]{GR25Classical}.
\end{proof}

\begin{proof}[
\Cref{prop:p031-scalar-openness-failure}]
Poineau--Turchetti consider an algebraically closed complete field $K$
which is not spherically complete, a maximally complete algebraically
closed extension $K'/K$, and a type--$4$ point of
$\mathbf A_K^{1,\mathrm{an}}$.  They identify its inverse image under
scalar projection as a positive-radius closed disc
\cite[Lemma~5.6.2]{PT21SchottkyI}; their
\cite[Remark~5.6.4]{PT21SchottkyI} records explicitly that the projection
is not open in general.  Thus the open-map hypothesis in
\Cref{lem:p031-open-pullback-usc} is not automatic, even for the
analytic affine line.  The final scope assertions follow because this
topological example does not supply a psh family for which the envelope
identity fails.
\end{proof}
